% ACTUAL CONTENT BEGINS ON LINE 500


\documentclass[12pt,a4paper]{article}

\newif\ifpublic

% Default to public for safety
\publictrue

% Allow override via command line flag
\ifdefined\PRIVATEBUILD
  \publicfalse
\fi

\ifpublic
  % Prevent accidental inclusion of private data
\ifdefined\PRIVATE
  \errmessage{ERROR: personal.tex included in PUBLIC build}
\fi

\newcommand{\StudentName}{REDACTED}
\newcommand{\CandidateNumber}{XXXX}
\newcommand{\CentreNumber}{XXXXX}
\newcommand{\SchoolName}{REDACTED}   % redacted
\else
  \input{personal.tex} % only for manual builds
\fi

% =============================
% PACKAGES
% =============================

% Allow UTF-8 encoded characters (e.g., accents, symbols)
\usepackage[utf8]{inputenc}

% T1 font encoding, improving hyphenation and font output
\usepackage[T1]{fontenc}

% page layout control (margins, paper size, etc.)
\usepackage{geometry}

% mathematical typesetting
\usepackage{amsmath}

% mathematical symbols (blackboard bold, etc.)
\usepackage{amssymb}

% page margins to 2.5 cm
\geometry{margin=2.5cm}

% Set headheight to accommodate fancyhdr
\setlength{\headheight}{13.6pt}

% Provides commands to control line spacing
\usepackage{setspace}

% Sets line spacing to 1.5 lines
\onehalfspacing

% Wrap text around figures
\usepackage{wrapfig}

% Font Awesome 5 for icons
\usepackage{fontawesome5}


% -----------------------------
% Graphics
% -----------------------------

% Loads graphicx in draft mode:
% - Images are replaced by empty boxes
% - Speeds up compilation during drafting
% Comment this out for final submission to show images
% \usepackage[draft]{graphicx}
\usepackage{graphicx}  % Uncomment this for final version with images

% Improves control over figure and table placement
\usepackage{float}

% professional table formatting
% (e.g., \toprule, \midrule, \bottomrule)
\usepackage{booktabs}
\usepackage{multirow}


% customisable numbering for lists
\usepackage{enumitem}

% clickable links for references, citations, URLs, and TOC
\usepackage{hyperref}

% smart quotation marks and language-aware quoting
\usepackage{csquotes}

% -----------------------------
% Bibliography
% -----------------------------

% Loads biblatex with:
% - biber as the backend
% - IEEE citation style
% - Maximum of 3 author names before "et al."
% - Back-references from bibliography to citations
\usepackage[
  backend=biber,
  style=ieee,
  maxbibnames=3,
  backref=true
]{biblatex}

% customisation of section titles
\usepackage{titlesec}

% bibliography db file location
\addbibresource{references.bib}

% smaller bibliography reference numbers
\DeclareFieldFormat{labelnumber}{\scriptsize#1}



% -----------------------------
% Styling and Layout Enhancements
% -----------------------------

% epigraphs (quotes at the start of chapters/sections)
\usepackage{epigraph}

% colours throughout the document
\usepackage{xcolor}
% Brand Color Definitions
\definecolor{brandBG}{HTML}{1A1A1A}
\definecolor{brandSurface}{HTML}{2D2D2D}
\definecolor{brandBorder}{HTML}{404040}
\definecolor{brandPrimary}{HTML}{6366F1}
\definecolor{brandSuccess}{HTML}{10B981}
\definecolor{brandWarning}{HTML}{F59E0B}
\definecolor{brandDanger}{HTML}{EF4444}
\definecolor{brandInfo}{HTML}{3B82F6}
\definecolor{brandTextPrimary}{HTML}{F3F4F6}
\definecolor{brandTextSecondary}{HTML}{9CA3AF}
\definecolor{brandTextDisabled}{HTML}{6B7280}

% Helper for color swatches
\newcommand{\swatch}[1]{\textcolor{#1}{\rule{1em}{1em}}}


% source code formatting and syntax highlighting
\usepackage{listings}

% vector graphics and diagrams
\usepackage{tikz}
\usetikzlibrary{arrows.meta, positioning, shapes.geometric, shapes.multipart, calc, bending}

% tree diagrams for hierarchical structures
\usepackage{forest}

% Proper Crow's Foot arrows declaration
\makeatletter
\pgfarrowsdeclare{crowsfoot}{crowsfoot}{
  \pgfarrowsleftextend{-.5\pgflinewidth}
  \pgfarrowsrightextend{-.5\pgflinewidth}
}{
  \pgfutil@tempdima=0.6pt
  \advance\pgfutil@tempdima by.4\pgflinewidth
  \pgfsetdash{}{+0pt}
  \pgfsetmiterjoin
  \pgfpathmoveto{\pgfqpoint{0pt}{-9\pgfutil@tempdima}}
  \pgfpathlineto{\pgfqpoint{-13\pgfutil@tempdima}{0pt}}
  \pgfpathlineto{\pgfqpoint{0pt}{9\pgfutil@tempdima}}
  \pgfpathmoveto{\pgfqpoint{-12pt}{-6pt}}
  \pgfpathlineto{\pgfqpoint{-12pt}{6pt}}
  \pgfusepathqstroke
}

\pgfarrowsdeclare{omany}{omany}{
  \pgfarrowsleftextend{-.5\pgflinewidth}
  \pgfarrowsrightextend{-.5\pgflinewidth}
}{
  \pgfutil@tempdima=0.6pt
  \advance\pgfutil@tempdima by.4\pgflinewidth
  \pgfsetdash{}{+0pt}
  \pgfsetmiterjoin
  \pgfpathmoveto{\pgfqpoint{0pt}{-9\pgfutil@tempdima}}
  \pgfpathlineto{\pgfqpoint{-13\pgfutil@tempdima}{0pt}}
  \pgfpathlineto{\pgfqpoint{0pt}{9\pgfutil@tempdima}}
  \pgfpathcircle{\pgfqpoint{-11.5pt}{0pt}}{3.5pt}
  \pgfusepathqstroke
}

\pgfarrowsdeclare{one}{one}{
  \pgfarrowsleftextend{-.5\pgflinewidth}
  \pgfarrowsrightextend{-.5\pgflinewidth}
}{
  \pgfutil@tempdima=0.6pt
  \advance\pgfutil@tempdima by.4\pgflinewidth
  \pgfsetdash{}{+0pt}
  \pgfsetmiterjoin
  \pgfpathmoveto{\pgfqpoint{-6pt}{-6pt}}
  \pgfpathlineto{\pgfqpoint{-6pt}{6pt}}
  \pgfusepathqstroke
}

\pgfarrowsdeclare{oone}{oone}{
  \pgfarrowsleftextend{-.5\pgflinewidth}
  \pgfarrowsrightextend{-.5\pgflinewidth}
}{
  \pgfutil@tempdima=0.6pt
  \advance\pgfutil@tempdima by.4\pgflinewidth
  \pgfsetdash{}{+0pt}
  \pgfsetroundjoin
  \pgfpathcircle{\pgfqpoint{0pt}{0pt}}{4\pgfutil@tempdima}
  \pgfusepathqstroke
}
\makeatother


% custom headers and footers
\usepackage{fancyhdr}


% -----------------------------
% Graphs and Charts
% -----------------------------

% plotting graphs (line charts, bar charts, etc.)
\usepackage{pgfplots}

% pgfplots v1.18 compatibility
\pgfplotsset{compat=1.18}

% pie charts
\usepackage{pgf-pie}



% =============================
% FORMATTING
% =============================

% customise headers and footers
% this is for the [School (Centre Number)] and [Candidate (Candidate Number)]
\pagestyle{fancy}

% clear all default header and footer fields
\fancyhf{}



% -----------------------------
% Header Content
% -----------------------------

% Left header:
% Displays the school name and centre number in small text
\lhead{\small \SchoolName (Centre Number: \CentreNumber)}

% Right header:
% Displays the candidate name and number in small text
% JEZLER - CHANGE THIS TO REAL NUMBER WHEN GIVEN
\rhead{\small \StudentName (Candidate Number: \CandidateNumber)}



% -----------------------------
% Footer Content
% -----------------------------

% Centre footer:
% Displays the current page number
\cfoot{\thepage}



% -----------------------------
% Header and Footer Lines
% -----------------------------

% thickness of the line under the header
\renewcommand{\headrulewidth}{0.4pt}

% Remove line above the footer
\renewcommand{\footrulewidth}{0pt}

% 1.5 line spacing throughout the document
\onehalfspacing



% -----------------------------
% PDF Metadata and Hyperlinks
% -----------------------------

% PDF properties and hyperlink appearance
\hypersetup{
    colorlinks=true,          % Enables coloured links instead of boxed links
    linkcolor=black,          % Colour of internal links (e.g. TOC)
    urlcolor=blue,            % Colour of URL links
    citecolor=blue,           % Colour of citation links
    pdflang=en-GB,            % Document language (British English)
    pdftitle={Computer Science NEA - REBORN}, % Title shown in PDF metadata
    pdfauthor={\StudentName [\CandidateNumber]}, % Author metadata
    pdfsubject={A-Level Computer Science Non-Exam Assessment}, % Subject metadata
    pdfkeywords={Computer Science, NEA, A-Level, Algorithms, Programming, Evaluation, Habit Tracker, Productivity, Self-improvement, AI}, % Search keywords
    plainpages=false         % Fixes issue with duplicate page numbers (e.g. i vs 1)
}



% -----------------------------
% Section Title Formatting
% -----------------------------

% Formats \section headings:
% - Normal font
% - Large size
% - Bold text
\titleformat{\section}
  {\normalfont\Large\bfseries} % Text style
  {\thesection}                % Section number
  {1em}                        % Space between number and title
  {}                           % Code before the title text


% Formats \subsection headings:
% - Slightly smaller than \section
% - Bold text
\titleformat{\subsection}
  {\normalfont\large\bfseries}
  {\thesubsection}
  {1em}
  {}




% -----------------------------
% Section Numbering Depth
% -----------------------------

% numbering down to paragraph level (1.1.1.1)
\setcounter{secnumdepth}{4}

% paragraph-level headings in the table of contents
\setcounter{tocdepth}{4}




% -----------------------------
% Paragraph as Subsubsubsection
% -----------------------------

% Redefines \paragraph to behave like a subsubsubsection
% Used for product analysis headings (e.g. Notion, Habitica, etc.)
\titleformat{\paragraph}
  {\normalfont\normalsize\bfseries}   % Normal size, bold
  {\thesubsubsection.\arabic{paragraph}} % Numbering like 1.1.1.1
  {1em}                              % Space between number and title
  {}

% Controls spacing around paragraph headings
\titlespacing*{\paragraph}
  {0pt}                              % Left indentation
  {3.25ex plus 1ex minus .2ex}       % Space before
  {1.5ex}                            % Space after




% -----------------------------
% Epigraph Formatting
% -----------------------------

% width of epigraph text
\setlength{\epigraphwidth}{0.4\textwidth}

% thickness of the rule beneath epigraphs
\setlength{\epigraphrule}{0.5pt}

% epigraph text to the right
\renewcommand{\textflush}{flushright}

% epigraph source to the right
\renewcommand{\sourceflush}{flushright}



% -----------------------------
% Tally Marks (TikZ)
% -----------------------------

% tally of five:
% Four vertical lines with a diagonal slash through them
\newcommand{\tallyfive}{%
  \begin{tikzpicture}[baseline=-3pt, scale=0.3]
    \draw[thick] (0,0) -- (0,1);
    \draw[thick] (0.3,0) -- (0.3,1);
    \draw[thick] (0.6,0) -- (0.6,1);
    \draw[thick] (0.9,0) -- (0.9,1);
    \draw[thick] (-0.2,0.2) -- (1.1,0.8);
  \end{tikzpicture}\hspace{2mm}%
}

% single tally line
% Can be repeated for 1–4 tallies
\newcommand{\tallyone}{%
  \begin{tikzpicture}[baseline=-3pt, scale=0.3]
    \draw[thick] (0,0) -- (0,1);
  \end{tikzpicture}\hspace{1mm}%
}



% =============================
% CODE FORMATTING
% =============================

% Global configuration for the listings package
\lstset{
    basicstyle=\ttfamily\small,   % Monospaced font, small size
    backgroundcolor=\color{gray!10}, % Light grey background
    frame=single,                 % Draws a border around code blocks
    breaklines=true,              % Automatically wraps long lines
    postbreak=\mbox{\textcolor{red}{$\hookrightarrow$}\space}, % Symbol indicating line break
    keywordstyle=\color{blue},    % Colour for language keywords
    commentstyle=\color{gray},    % Colour for comments
}

% Define a premium pseudocode style
\lstdefinestyle{pseudocode}{
    basicstyle=\ttfamily\small,
    keywordstyle=\color{blue!80!black}\bfseries,
    commentstyle=\color{green!50!black}\itshape,
    stringstyle=\color{red!80!black},
    numbers=left,
    numberstyle=\tiny\color{gray},
    stepnumber=1,
    numbersep=10pt,
    backgroundcolor=\color{gray!5},
    frame=single,
    rulecolor=\color{black!20},
    tabsize=4,
    captionpos=b,
    breaklines=true,
    breakatwhitespace=true,
    showspaces=false,
    showtabs=false,
    xleftmargin=2.5em,
    framexleftmargin=1.5em,
    morekeywords={FUNCTION, IF, THEN, ELSE, RETURN, FOR, EACH, IN, DO, END, DELETE, TRUE, FALSE, LENGTH},
    morecomment=[l]{//},
    literate={<-}{{$\leftarrow$}}1 {<=}{{$\leq$}}1 {>=}{{$\geq$}}1,
}

\usepackage{algorithm}
\usepackage{algpseudocode}

% Define AQA-style alias commands if strictly needed, 
% but standard algorithmicx usually suffices.
\renewcommand{\algorithmicrequire}{\textbf{INPUT:}}
\renewcommand{\algorithmicensure}{\textbf{OUTPUT:}}


\usepackage{minted}
\setminted{
  fontsize=\small,
  frame=lines,
  framesep=6pt,
  linenos,
  breaklines,
  tabsize=2,
  bgcolor=gray!5
}


% =============================
% DOCUMENT
% =============================
\begin{document}
% ---------- TITLE PAGE ----------
 \hypersetup{pageanchor=false}
 \begin{titlepage}
    \centering
    \vspace*{3cm}
    {\Huge \textbf{Computer Science NEA}}\\[0.5cm]
    {\Large REBORN --- An adaptive Habit Tracker}\\[1cm]

    {\Large \StudentName}\\[0.3cm]
    {\large Candidate Number: \textbf{\CandidateNumber} \quad | \quad Centre Number: \textbf{\CentreNumber}}\\[0.3cm]
    {\large Centre Name: \textbf{\SchoolName}}\\[0.3cm]
    {\large Qualification Code: 7517}\\[0.3cm]
    {\large \today}\\[3cm]

    \rule{0.9\textwidth}{0.4pt}\\[0.3cm]
    \begin{minipage}{0.8\textwidth}
    \centering
    \emph{“Civilization advances by extending the number of operations we can perform without thinking about them.” --- A. North Whitehead}
    \end{minipage}\\[0.3cm]
    \rule{0.9\textwidth}{0.4pt}
\end{titlepage}
\hypersetup{pageanchor=true}


\pagenumbering{roman}
\tableofcontents
\listoffigures
\listoftables
\clearpage
\pagenumbering{arabic}


% ================================================================================
% ------------------------------- Introduction -----------------------------------
% ================================================================================
\section{Introduction}

\epigraph{\enquote{Make it so easy you can't say no.}}{\textit{James Clear, Atomic Habits} \cite{clear2018atomic}}

Modern digital systems increasingly compete for human attention through algorithmic optimisation, yet provide little computational support for helping users regulate behaviour and build sustainable habits. Many widely used platforms prioritise engagement over user agency, contributing to inconsistent routines and behavioural instability.

This project, \textbf{REBORN}, aims to build an intelligent habit-tracking system that automates scheduling, analyses behavioural risk, and dynamically adapts difficulty. When software handles the heavy planning, users can focus their energy on execution.

% ---------- 2. ANALYSIS ----------
% Analysis (7-9 marks)
% "Fully scoped analysis of a real problem, presented in a way that a third party can understand. Requirements fully documented in a set of measurable and appropriate specific objectives, covering all required functionality of the solution or areas of investigation. Requirements arrived at by considering, through dialogue, the needs of the intended users of the system, or recipients of the outcomes for investigative projects. Problem sufficiently well modelled to be of use in subsequent stages."

\clearpage
% ===================================================================================
\section{Analysis}
% ===================================================================================

\subsection{Problem Definition}

% Describing how manual habit tracking is exhausting and current technology is not intelligent enough to reduce this problem.
Many people experience \emph{self-regulation failure} - they want to improve, but the cognitive effort needed to plan and continuously maintain habits often leads to burnout \cite{baumeister1996self, baumeister1998ego}. Technology is not only the largest source of procrastination, but also doesn't provide intelligent ways to reduce this problem \cite{sirois2016procrastination}. Habit formation, therefore, is both a computational challenge and a psychological one. Tools that can automate planning, predict high-risk periods, and adapt interventions in real time are therefore necessary.


\subsection{Users}
% Describing the primary users of the system and their needs.
% This includes: 
%  A level students
%  Professionals
%  Fitness enthusiasts
%  General users who want to improve their habits
The primary users of REBORN are students, especially those studying multiple A-Levels and struggling with procrastination and the mental load of planning habits. For this group, automation and predictive intervention allow them to focus on doing the habits instead of organisation.  

\vspace{0.5em}
The system is also intended for a larger audience, such as:
\begin{itemize}
    \item Professionals who need structured routines and work-life balance.
    \item Individuals maintaining regular fitness or health-related habits.
    \item Users interested in general self-improvement and long-term behaviour change.
\end{itemize}

All of these user groups experience forms of self-regulation failure, where manual habit tracking can lead to cognitive overload and burnout.
% ================================================================================
% --------------------- Existing Solutions - Secondary Research ------------------
% ================================================================================
\newpage
\subsection{Secondary Research}

% The current state of habit tracking apps, and why we need REBORN.
Current habit-tracking and productivity applications differ widely in how much computation they perform. Some are manual-input database systems, others are game-like behavioural frameworks. An analysis of these tools shows that most apps use psychological principles (loss aversion/variable reward schedules) instead of algorithmic reasoning.

To justify the development of REBORN, I will contrast existing tools with a theoretical \enquote{gold standard} of computational habit formation, and then find specific technical gaps in scheduling, prediction, and how systems adapt to user context over time.

\subsubsection{PCS Framework}
% Some theory on PCS (the strongest framework for habit formation).
The strongest computational framework for habit formation is the \textbf{Predicting Context Sensitivity (PCS)} method \cite{buyalskaya2023what}. Unlike the \emph{21-day myth}, this machine-learning methodology models habit formation as a continuous measure of context-sensitive predictability.

% Explain what LASSO is here (machine learning tool used in habit research)
Using \textbf{LASSO} (Least Absolute Shrinkage and Selection Operator) regression---a regularised machine-learning method that identifies the most predictive contextual variables while reducing overfitting---on datasets with 40 million+ observations, Buyalskaya et al.\cite{buyalskaya2023what} showed that habit formation timelines are domain-specific. The estimated times range from a few weeks for simple hygiene tasks to 2--3 months for more complex behaviour like regularly going to the gym.

% Why REBORN is needed
Most importantly, no current app makes use of this kind of context data (e.g., time lags, or day-of-week streaks) to predict user behaviour.

% ----------------------------------------------------
% ---------- ANALYSIS OF EXISTING SOLUTIONS ----------
% ----------------------------------------------------

\subsubsection{Existing Systems}
This section analyses the four most popular habit tracking and task management apps: Notion, Habitica, Streaks, and Todoist. While other applications exist, these four represent the distinct architectural archetypes currently dominating the market.

% ---------------------------------------
% 1. DESCRIPTIONS AND FUNCTIONALITY
% ---------------------------------------
\newpage
\paragraph{Description and Functionality}

\textbf{Notion} \\
Notion is a workspace that acts like a relational database wrapper. The user defines schemas (tables) and views (Kanban, Calendar) manually. The backend uses a SQL-like structure which is then served to a React frontend.
\begin{figure}[H]
    \centering
    \includegraphics[width=0.8\textwidth, trim=0 50 0 50, clip]{images/notion_secondary_research.png} 
    \caption{A Notion Habit Tracker showing the manual checkbox system.}
    \label{fig:notion}
\end{figure}

\noindent\textbf{Habitica} \\
Habitica is a gamified task manager that uses operant conditioning \cite{skinner1938}. It wraps a CRUD system in RPG mechanics, where doing tasks gives \emph{Gold/XP} and missing them causes \emph{Health Loss}.
\begin{figure}[H]
    \centering
    \includegraphics[width=0.8\textwidth]{images/habitica_secondary_research.png} 
    \caption{The Habitica dashboard showing the RPG avatar and health bars.}
    \label{fig:habitica}
\end{figure}

\noindent\textbf{Streaks} \\
Streaks is a minimalist habit tracker integrated into the Apple ecosystem. Unlike database tools, it focuses entirely on the \emph{Don't Break the Chain} psychological method \cite{clear2018atomic}. Technically, it is notable for its API integration with Apple HealthKit \cite{apple2025healthkit}, allowing it to automatically mark habits (e.g., \enquote{Walk 5,000 steps}) as complete without user input.
\begin{figure}[H]
    \centering
    \includegraphics[width=0.8\textwidth]{images/streaks_secondary_research.jpeg} 
    \caption{The Streaks interface showing its main features.}
    \label{fig:streaks}
\end{figure}

\noindent\textbf{Todoist} \\
Todoist is a task management system, designed for one-off tasks. In earlier versions, Todoist experimented with automated rescheduling through its \textit{Smart Schedule} feature, which attempted to suggest new due dates for overdue tasks. Although it was discontinued because of accuracy issues \cite{todoist2020smartschedule}, third-party integrations like Trevor AI \cite{trevorai2025} can provide similar functionality, which apply heuristic time-blocking and NLP to parse tasks and suggest times for them based on available calendar gaps.
\begin{figure}[H]
    \centering
    \includegraphics[width=0.8\textwidth, trim=0 100 0 150, clip]{images/todoist_secondary_research.jpg} 
    \caption{Todoist's ``Smart Schedule'' interface suggesting times for tasks \cite{todoist_smart_schedule_image_2016}.}
    \label{fig:todoist}
\end{figure}

% ---------------------------------------
% 2. ADVANTAGES AND DISADVANTAGES
% ---------------------------------------
\newpage
\paragraph{Analysis of Advantages and Disadvantages}

The following tables breakdown the strengths and weaknesses of each system's architecture and user experience model.

\vspace{2em}
\noindent\textbf{Notion}
\begin{table}[H]
    \centering
    \begin{tabular}{|p{0.45\linewidth}|p{0.45\linewidth}|}
    \hline
    \textbf{Advantages} & \textbf{Disadvantages} \\
    \hline
    Flexible data modeling that allows complex relationship definitions. & \textbf{The \emph{Blank Slate} problem:} Users have to build the system themselves; there is no built-in logic. (without templates)\\
    \hline
    Visual customisation reduces initial friction for design-oriented users. & \textbf{No constraint satisfaction:} It cannot detect schedule conflicts or optimise time slots automatically \cite{russell2020ai}. \\
    \hline
    Excellent database management for historical record keeping. & Entirely manual input leads to high administrative load (cognitive friction). \\
    \hline
    \end{tabular}
\end{table}

\noindent\textbf{Streaks}
\begin{table}[H]
    \centering
    \begin{tabular}{|p{0.45\linewidth}|p{0.45\linewidth}|}
    \hline
    \textbf{Advantages} & \textbf{Disadvantages} \\
    \hline
    \textbf{API Automation:} Integration with Apple HealthKit removes \emph{administrative friction} by tracking physical habits automatically. & \textbf{\emph{All-or-Nothing}:} It uses binary logic. If the user misses one day, the streak resets to zero. This often causes the \emph{What-the-Hell Effect} \cite{polivy1975perceived}, where users abandon the app after a failure. \\
    \hline
    \textbf{Visual Feedback:} The circular UI provides visual status of daily progress, using \emph{loss aversion} for engagement \cite{kahneman2011thinking}. & \textbf{No Scheduling Logic:} The app tracks \textit{completion}, but not \textit{allocation}. It cannot help the user find time in a busy schedule to actually perform the habit. \\
    \hline
    Minimalist design reduces cognitive load compared to complex tools like Notion. & \textbf{No adaptive difficulty:} The target remains static regardless of user burnout. \\
    \hline
    \end{tabular}
\end{table}

\newpage
\noindent\textbf{Habitica}
\begin{table}[H]
    \centering
    \begin{tabular}{|p{0.45\linewidth}|p{0.45\linewidth}|}
    \hline
    \textbf{Advantages} & \textbf{Disadvantages} \\
    \hline
    Gamification (XP/Gold) uses dopamine loops \cite{eyal2014hooked} to increase engagement. & \textbf{Static Recurrence:} It relies on fixed rules (e.g., \enquote{Every Mon/Tue}) instead of adaptive scheduling. \\
    \hline
    Social accountability features (parties/guilds) help with external motivation. & \textbf{Reactive Punishment:} It creates a \emph{death spiral}. If a user is busy and misses tasks, the game punishes them, demotivating them further. \\
    \hline
    Visualisation of progress from avatar customisation and item unlocks shows long-term habit progress. & No awareness of external schedules or calendar conflicts. \\
    \hline
    \end{tabular}
\end{table}

\noindent\textbf{Todoist}
\begin{table}[H]
    \centering
    \begin{tabular}{|p{0.45\linewidth}|p{0.45\linewidth}|}
    \hline
    \textbf{Advantages} & \textbf{Disadvantages} \\
    \hline
    NLP (Natural Language Processing) allows for quick task entry using natural language \cite{todoistQuickAdd}. & \textbf{Task vs. Habit:} The system focuses on one-off tasks and does not have a risk or failure model for long-term habit formation. \\
    \hline
    Heuristic time-blocking reduces the cognitive load of manual scheduling. & Scheduling behaviour is static; repeated failure does not reduce workload or adjust task frequency. \\
    \hline
    Deep ecosystem integrations (email, calendar, browser, automation tools) allow Todoist to function as a centralised task capture hub across contexts and devices. & Closed-source algorithms prevent inspection or tuning for individual habit difficulty curves. \\
    \hline
    \end{tabular}
\end{table}


% ---------------------------------------
% 3. IMPLICATIONS
% ---------------------------------------
\paragraph{Implications for Proposed System}

By investigating these systems, I have identified specific features to adopt and flaws to improve upon in the proposed solution.

\newpage
\noindent\textbf{From Notion:} \\
While a relational database is important for storing habit history, relying on the user to manually schedule every event will likely lead to burnout.
\begin{itemize}
    \item \textbf{I will include:} A strong relational database backend similar to Notion's architecture for data integrity.
    \item \textbf{I will improve:} Instead of manual entry, I will implement a Constraint Satisfaction Scheduler to automate the data entry process, solving the \emph{Blank Slate} problem.
\end{itemize}

\vspace{0.5em}
\noindent\textbf{From Habitica:} \\
Feedback mechanisms are effective, but \enquote{punishment} (losing HP) is flawed if the user is busy or overwhelmed.
\begin{itemize}
    \item \textbf{I will include:} A feedback loop system, but instead of RPG stats, I will use a \textbf{Simulation} to visualise habit stability.
    \item \textbf{I will improve:} The system will be predictive rather than reactive. By using \textbf{Logistic Regression}, my system will warn the user of risk \textit{before} they fail, instead of punishing them \textit{after}.
\end{itemize}

\vspace{0.5em}
\noindent\textbf{From Streaks:} \\
Visualising progress is very important, but the \emph{streak counter} model is too fragile for students with variable workloads.
\begin{itemize}
    \item \textbf{I will include:} A strong visual representation of habit health, similar to the \emph{rings} in Streaks, to provide immediate feedback.
    \item \textbf{I will improve:} Instead of a binary \emph{streak} that breaks (0 or 1), I will use my \textbf{simulation}. This treats habit consistency as a continuous variable (stability). If a user misses a day, the stability decays instead of resetting to zero, which is mathematically more forgiving and encourages the user to recover.
\end{itemize}

\vspace{0.5em}
\noindent\textbf{From Todoist:} \\
Algorithmic scheduling is desirable, but fixed heuristics fail when psychological factors and user failure are not modelled.
\begin{itemize}
    \item \textbf{I will include:} Automated time-blocking to ensure habits and tasks are placed within realistic daily constraints.
    \item \textbf{I will improve:} The use of a \emph{Spaced Repetition} algorithm to dynamically reduce habit frequency when users struggle, preventing overload and burnout.
    \item \textbf{I will avoid:} Set heuristic scheduling, by integrating adaptive control and predictive modelling to maintain long-term habit.
\end{itemize}


% ---------------------------------------
% 4. TECHNICAL GAPS
% ---------------------------------------
\newpage
\subsubsection{Identification of Technical Gaps}

The comparison of existing tools shows three algorithmic gaps in the current market that \textbf{REBORN} aims to solve:

\begin{enumerate}
    \item \textbf{Constraint Satisfaction Solvers (CSP):} 
    Current apps rely on manual time-picking. No system treats the daily schedule as a combinatorial search problem that can be solved using backtracking or constraint propagation to guarantee conflict-free allocation \cite{russell2020ai}.
    
    \item \textbf{Predictive Logistic Regression:} 
    Applications are reactive; they show historical graphs instead of forecasting. None use the PCS framework's approach (Logistic Regression) to calculate a probability $P(fail)$ for the upcoming day to trigger preventative interventions.
    
    \item \textbf{Static Frequency vs. Adaptive Control:} 
    Habit trackers use static schedules (e.g., ``Daily''). They do not implement control loops (like the SM-2 spaced repetition algorithm \cite{wozniak1990sm2}) to adjust frequency based on user performance, for example, by scaling back load to prevent burnout.
\end{enumerate}

\begin{table}[H]
\centering
\resizebox{\textwidth}{!}{%
\begin{tabular}{|l|c|c|c|c|}
\hline
\textbf{System} & \textbf{Scheduling Logic} & \textbf{Predictive Model} & \textbf{Adaptivity} & \textbf{Primary Domain} \\ \hline
Notion & Manual (User defined) & None & None & Database \\ \hline
Habitica & Fixed Recurrence & None & None & Gamification \\ \hline
Streaks & Fixed Recurrence & None & None & Health Tracking \\ \hline
Todoist + Trevor & Heuristic Time-blocking & Duration Estimation & None & Task Management \\ \hline
\textbf{REBORN} & \textbf{Constraint Solver (DFS)} & \textbf{Logistic Regression} & \textbf{SM-2 Algorithm} & \textbf{Habit Optimisation} \\ \hline
\end{tabular}%
}
\caption{Comparison of existing solutions against the proposed system.}
\label{tab:comparison}
\end{table}

% ---------- CONCLUSION ----------
\subsubsection{Conclusion}
The existing market is saturated with tools that record data but can't compute solutions. Todoist already shows how algorithmic scheduling can help with one-off tasks and academic papers on the PCS framework \cite{buyalskaya2023what} led mathematically with prediction. No system integrates these into a single application. 



% End of new appendices

\newpage

% ================================================================================
% ------------------------------------- End --------------------------------------
% ================================================================================
\clearpage
\thispagestyle{empty}
\vspace*{\fill}

\begin{center}
\rule{0.6\textwidth}{0.4pt}\\[0.4cm]
\textit{REBORN --- Computer Science NEA}\\[0.3cm]
\rule{0.6\textwidth}{0.4pt}
\end{center}

\vspace*{\fill}


% ================================================================================
% ---------------------------------- References ----------------------------------
% ================================================================================
\clearpage
\addcontentsline{toc}{section}{References}
\printbibliography[title=References]

\end{document}